\chapter{Reproduction and Study Protocol}\label{app:repro}

\section{Reproducing the results}

Every result in Chapter~\ref{ch:results} is produced by a script in \code{scripts/research/}. Each
script writes a JSON file to \code{eval/research/}, which the report builder then reads. From the
repository root, in the project's virtual environment:

\begin{verbatim}
python -m scripts.research.validate_labeller --seed 11 --tag _heldout
python -m scripts.research.run_simulation --runs 200
python -m scripts.research.run_simulation --runs 200 --only E6,E7
python -m scripts.research.evaluate_difficulty_models
python -m scripts.research.fetch_cohort --per-band 60
python -m scripts.research.analyze_cohort --workers 14 --max-games 60
python -m scripts.research.evaluate_weakness_models --write-norms
python -m scripts.research.classify_rating_band
python -m scripts.research.make_figures
python -m scripts.research.build_report
pytest tests/
\end{verbatim}

\noindent Cohort collection needs network access and takes roughly an hour. Cohort analysis takes
about ten hours on the development laptop. Every other step runs in minutes.

\section{Reproducing PuzzleNet}

The neural pipeline lives in \code{scripts/neural/} and writes to \code{eval/neural/}:

\begin{verbatim}
python -m scripts.neural.build_dataset --workers 12
python -m scripts.neural.run_experiments
python -m scripts.neural.engine_pv_check --n 2000 --workers 6
python -m scripts.neural.evaluate_puzzlenet
python -m scripts.neural.run_label_simulation --runs 200
python -m scripts.neural.relabel_cohort --pvs --workers 8
python -m scripts.neural.relabel_cohort
python -m scripts.research.evaluate_weakness_models
    --analysis-dir data/research/cohort/analysis_puzzlenet
    --out-name cohort_evaluation_puzzlenet.json
python -m scripts.research.evaluate_weakness_models
    --analysis-dir data/research/cohort/analysis_rules_rerun
    --out-name cohort_evaluation_rules_rerun.json
python -m scripts.research.evaluate_difficulty_models
python -m scripts.neural.make_figures
\end{verbatim}

\noindent Encoding the puzzle database takes about 50 minutes and produces 2.4\,GB of features in
\code{data/neural/} (not in version control). Training the production model takes about half an hour
on a laptop CPU; the ablations and the learning curve add roughly two hours. Re-running Stockfish on
the cohort's 40{,}067 critical positions takes about an hour. The trained model itself,
\code{src/data/models/puzzlenet.npz} (3.6\,MB), \emph{is} in version control, so the evaluation
scripts can be run without retraining. The two \code{evaluate\_weakness\_models} commands are single
lines; they are wrapped here for the page width.

\section{Summary of the pre-registered human study}

The full protocol is in \code{docs/USER\_STUDY\_PROTOCOL.md}.
\begin{itemize}
  \item \textbf{Design.} A within-subject crossover. Each participant trains for two weeks under the
        adaptive policy and two weeks under the control (random category, Elo $\pm 300$ band), with the
        order counterbalanced and 20 puzzles a day.
  \item \textbf{Primary outcome.} Change in score on matched probe sets that are never served during
        training.
  \item \textbf{Secondary outcome.} Change in the opportunity-normalised miss rate in the participant's
        real games.
  \item \textbf{Analysis.} A linear mixed model of condition by period, with a random intercept per
        participant.
  \item \textbf{Sample size.} 34 paired participants for $d=0.5$ at 80\% power (exact noncentral-$t$
        calculation); 40 to allow for attrition.
  \item \textbf{Ethics.} Approval is required before recruitment. \authornote{status}
\end{itemize}
